%%% Exemplo de utilização da classe ITA
%%%
%%%   por        Fábio Fagundes Silveira   -  ffs [at] ita [dot] br
%%%              Benedito C. O. Maciel     -  bcmaciel [at] ita [dot] br
%%%              Giovani Volnei Meinertz   -  giovani [at] ita [dot] br
%%%    	         Hudson Alberto Bode       -  bode [at] ita [dot]br
%%%    	         P. I. Braga de Queiroz    -  pi [at] ita [dot] br
%%%    	         Jorge A. B. Gripp         -  gripp [at] ita [dot] br
%%%    	         Juliano Monte-Mor         -  jamontemor [at] yahoo [dot] com [dot] br
%%%    	         Tarcisio A. B. Gripp      -  tarcisio.gripp [at] gmail [dot] com
%%%
%%%   Versão para overleaf:
%%%   por        Alejandro A. Rios Cruz    - aarc.88@gmail.com
%%%              Saulo Gómez               - sagomezs@unal.edu.co
%%%              Ocimar Santos             - ocimar.acad@gmail.com
%%%
%%%   Template disponibilizado em:
%%%              Overleaf: https://pt.overleaf.com/latex/templates/thesis-template-aeronautics-institute-of-technology-ita/yhfrqqydpygk
%%%
%%%   Contribuia você também!
%%%              GitHub:   https://github.com/AlejandroRios/Template_Thesis_ITA
%%%
%%%  IMPORTANTE: O texto contido neste exemplo nao significa absolutamente nada.  :-)
%%%              O intuito aqui eh demonstrar os comandos criados na classe e suas
%%%              respectivas utilizacoes.
%%%
%%%  Tese.tex  2016-08-25
%%%  $HeadURL: http://www.apgita.org.br/apgita/teses-e-latex.php $
%%%
%%% ITALUS
%%% Instituto Tecnológico de Aeronáutica --- ITA, Sao Jose dos Campos, Brasil
%%%                   http://groups.yahoo.com/group/italus/
%%% Discussion list: italus {at} yahoogroups.com
%%%
%++++++++++++++++++++++++++++++++++++++++++++++++++++++++++++++++++++++++++++++
% Para alterar o TIPO DE DOCUMENTO, preencher a linha abaixo \documentclass[?]{?}
%   \documentclass[tg]{ita}			= Trabalho de Graduacao
%   \documentclass[tgfem]{ita}	= Para Engenheiras
%   								msc     		= Dissertacao de Mestrado
%   								mscfem   		= Para Mestras
%   								dsc      		= Tese de Doutorado
%   								dscfem   		= Para Doutoras
%   								quali    		= Exame de Qualificacao
%   								qualifem 		= Exame de Qualificacao para Doutoras
% Para 'Draft Version'/'Versao Preliminar' com data no rodape, adicionar 'dv':
%   \documentclass[dsc, dv]{ita}
% Para trabalhos em Inglês, adicionar 'eng':
%   \documentclass[dsc, eng]{ita}
%		\documentclass[dsc, eng, dv]{ita}
%++++++++++++++++++++++++++++++++++++++++++++++++++++++++++++++++++++++++++++++
\documentclass[tg]{ita}    % ITA.cls based on standard book.cls
% Quando alterar a classe, por exemplo de [msc] para [msc, eng]) rode mais uma vez o botão BUILD OUTPUT caso haja erro
\usepackage{ae}
\usepackage{graphicx}
\usepackage{epsfig}
\usepackage{amsmath}
\usepackage{amssymb}
\usepackage{subfig}
\usepackage{multirow}
\usepackage{float}
\usepackage{amsthm}
\usepackage{url}         % formats URL addresses properly
\usepackage{appendix}    % allows appendix section to be included
\usepackage{lscape}      % allows a page to be rendered in landscape mode
\usepackage{multicol}    % allows text in multi columns
\usepackage{cancel}      % needed to show canceled terms in equations
\usepackage{lettrine}
\usepackage{float}
\usepackage{placeins}


%HHHHHHHHHHHHHHHHHHHHHHHHHHHHHHHHHHHHHHHHHHHHHHHHHHHHHHHHHHHHHHHHHHHHHHHHHHHHHHHHHHHHHHHHHHHHHHHHHHHHHHHHHHHH
%\usepackage{subfigure}
%\usepackage{subfigmat}
%PACOTEFIGURAS_SE _ERRADO_ESXCLUIR_ACIMA
\usepackage{booktabs}
%PACOTETABELAS_SE _ERRADO_ESXCLUIR_ACIMA
%HHHHHHHHHHHHHHHHHHHHHHHHHHHHHHHHHHHHHHHHHHHHHHHHHHHHHHHHHHHHHHHHHHHHHHHHHHHHHHHHHHHHHHHHHHHHHHHHHHHHHHHHHHHH

%++++++++++++++++++++++++++++++++++++++++++++++++++++++++++++++++++++++++++++++
% Espaçamento padrão de todo o documento
%++++++++++++++++++++++++++++++++++++++++++++++++++++++++++++++++++++++++++++++
\onehalfspacing

%singlespacing Para um espaçamento simples
%onehalfspacing Para um espaçamento de 1,5
%doublespacing Para um espaçamento duplo

%++++++++++++++++++++++++++++++++++++++++++++++++++++++++++++++++++++++++++++++
% Identificacoes (se o trabalho for em inglês, insira os dados em inglês)
% Para entradas abreviadas de Professora (Profa.) em português escreva: Prof$^\textnormal{a}$.
%++++++++++++++++++++++++++++++++++++++++++++++++++++++++++++++++++++++++++++++
\course{Engenheria Aeroespacial}  % Programa de PG ou Curso de Graduação
%\area{Aircraft Design} % Área de concentração na PG (Não utilizado no caso de TG)

% Autor do trabalho: Nome Sobrenome
\authorgender{masc}                     %sexo: masc ou fem
\author{Pedro Kuntz}{Puglia}
\itaauthoraddress{Rua H8C, Ap. 303}{12.228- 462}{São José dos Campos- SP}

% Titulo da Tese/Dissertação
\title{Caracterização de sistema de propulsão a gás frio com empuxo vetorial}

% Orientador
\advisorgender{masc}                    % masc ou fem
\advisor{Prof.~Dr.}{Leonardo Gouvêa}{ITA}

% Coorientador (Caso não haja coorientador, colocar ambas as variáveis \coadvisorgender e \coadvisor comentadas, com um % na frente)
% \coadvisorgender{fem}									% masc ou fem
% \coadvisor{Prof$^\textnormal{a}$.~Dr$^\textnormal{a}$.}{Doralice Serra}{OVNI}

% Pró-reitor da Pós-graduação
% \bossgender{masc}												% masc ou fem
% \boss{Prof.~Dr.}{John von Neumann}

%Coordenador do curso no caso de TG
\bosscoursegender{fem}									% masc ou fem
\bosscourse{Profa.~Dra.}{Cristiane Martins}

% Palavras-Chaves informadas pela Biblioteca -> utilizada na CIP
\kwcip{Propulsão}
\kwcip{Empuxo Vetorial}
\kwcip{Gás Frio}

% membros da banca examinadora

% \examiner{Prof. Dr.}{Alan Turing}{Presidente}{ITA}
% \examiner{Prof. Dr.}{Linus Torwald}{}{UXXX}
% \examiner{Prof. Dr.}{Richard Stallman}{}{UYYY}
% \examiner{Prof. Dr.}{Donald Duck}{}{DYSNEY}
% \examiner{Prof. Dr.}{Mickey Mouse}{}{DISNEY}

% Data da defesa (mês em maiúsculo, se trabalho em inglês, e minúsculo se trabalho em português)
\date{23}{novembro}{2022}

% Número CDU - (somente para TG)
\cdu{???.??}

% Glossario
\makeglossary
\frontmatter

\begin{document}
%%%%%%%%%%%%%%%%%%%%%%%%%%%%%%%%%%%%%%%%%
%%% - - Modelo para relatório ITA - - %%%
%%%%%%%%%%%%%%%%%%%%%%%%%%%%%%%%%%%%%%%%%
%%%    Documente as mudanças e os     %%%
%%%       upgrades na estrutura       %%%
%%%%%%%%%%%%%%%%%%%%%%%%%%%%%%%%%%%%%%%%%

\documentclass[12pt, a4paper]{article}

% Fonte e símbolos
% \usepackage{fontspec}
% \setmainfont{Times New Roman}
\usepackage[brazilian]{babel}
% Tamanho padrão da folha
\usepackage[left=1.5cm,right=1.5cm,top=1.5cm,bottom=1.5cm]{geometry}
\usepackage{csquotes}
\usepackage{latexsym}
\usepackage{amsmath}
\usepackage{textcomp}
\usepackage{gensymb}
\usepackage{indentfirst}
\usepackage{amssymb}
\usepackage{graphicx}
\usepackage{pbox}
\usepackage[dvipsnames]{xcolor}
\usepackage{float}
\usepackage{verbatim}
\usepackage{adjustbox}


% Top, mid e bottomrule para tabelas
\usepackage{booktabs}
\usepackage{multirow}

% Gerenciamento de referências
% É possível adicionar "style=abnt"
\usepackage[
backend=biber,
sorting=none,
style=abnt-numeric
]{biblatex}
\addbibresource{references.bib}

% Tamanho 12 para os títulos de seção e subsseção
\usepackage[tiny]{titlesec}

% Package para fazer referências no corpo do texto
\usepackage{hyperref}

% Packages para uso de subfigure
\usepackage{caption}
\usepackage{subcaption}
\usepackage{ragged2e}

\begin{document}

\begin{titlepage}
\begin{center}

\begin{figure}
\begin{center}
\includegraphics[width=9cm]{Images/ita-logo.png}
\end{center}
\end{figure}

\begin{huge}
\bf{\sc{Instituto Tecnológico de Aeronáutica}} \\
\vspace*{0.15in}
\end{huge}

\begin{Large}
\bf{\sc{Divisão de Engenharia Aeronáutica e Aeroespacial}}\\
\end{Large}

\begin{Large}
\vspace*{0.2in}
\textbf{Departamento de Eletrônica (IEE)} \\
\textbf{Projeto de ELE-27} \\

\vfill

\textbf{CubeSat}
\end{Large}

\vfill


\begin{large}
    \begin{FlushLeft}
        \textbf{Professor} \\
        Tertuliano Ribeiro Pinto \\
    \end{FlushLeft}
\end{large}

\begin{large}

\begin{FlushLeft}
    \textbf{Alunos} \\
    Ana Beatriz da Silva Machado\\
    Caio Jansen Accioly\\
    Lucas Ramalho Rocha\\
    Luísa Dallmann\\
    Mahmud Mohamad Hussein Ali Neto\\
    Mateus Silva Borges\\
    Natan Campistano de Mello\\
    Pedro Kuntz Puglia \\
\end{FlushLeft}

\end{large}
\vspace*{0.1in}
\Large{2023}
\end{center}
\end{titlepage}




\input{Tex/0.Sumário.tex}

\input{Tex/9.Referências.tex}

\end{document}

\mainmatter
% Os capitulos comecam aqui
\chapter{Introdução}
\section{Motivação e objetivos}




% REFERENCIAS BIBLIOGRAFICAS
\renewcommand\bibname{\itareferencesnamebabel} %renomear título do capítulo referências
\bibliography{Referencias/referencias}

% Apendices
\appendix

% Anexos
\annex

% Glossario
%\itaglossary
%\printglossary

% Folha de Registro do Documento
% Valores dos campos do formulario
\FRDitadata{25 de março de 2015}
\FRDitadocnro{DCTA/ITA/DM-018/2015} %(o número de registro você solicita a biblioteca)
\FRDitaorgaointerno{Instituto Tecnológico de Aeronáutica -- ITA}
%Exemplo no caso de pós-graduação: Instituto Tecnol{\'o}gico de Aeron{\'a}utica -- ITA
\FRDitapalavrasautor{Cupim; Cimento; Estruturas}
\FRDitapalavrasresult{Cupim; Dilema; Construção}
%Exemplo no caso de graduação (TG):
%\FRDitapalavraapresentacao{Trabalho de Graduação, ITA, São José dos Campos, 2015. \NumPenultimaPagina\ páginas.}
%Exemplo no caso de pós-graduação (msc, dsc):
\FRDitapalavraapresentacao{ITA, São José dos Campos. Curso de Mestrado. Programa de Pós-Graduação em Engenharia Aeronáutica e Mecânica. Área de Sistemas Aeroespaciais e Mecatrônica. Orientador: Prof.~Dr. Adalberto Santos Dupont. Coorientadora: Prof$^\textnormal{a}$.~Dr$^\textnormal{a}$. Doralice Serra. Defesa em 05/03/2015. Publicada em 25/03/2015.}
\FRDitaresumo{Aqui começa o resumo do referido trabalho. Não tenho a menor idéia do que colocar aqui. Sendo assim, vou inventar. Lá vai: Este trabalho apresenta uma metodologia de controle de posição das juntas passivas de um manipulador subatuado de uma maneira subótima. O termo subatuado se refere ao fato de que nem todas as juntas ou graus de liberdade do sistema são equipados com atuadores, o que ocorre na prática devido a falhas ou como resultado de projeto. As juntas passivas de manipuladores desse tipo são indiretamente controladas pelo movimento das juntas ativas usando as características de acoplamento da dinâmica de manipuladores. A utilização de redundância de atuação das juntas ativas permite a minimização de alguns critérios, como consumo de energia, por exemplo.
Apesar da estrutura cinemática de manipuladores subatuados ser idêntica a do totalmente atuado, em geral suas caraterísticas dinâmicas diferem devido a presença de juntas passivas. Assim, apresentamos a modelagem dinâmica de um manipulador subatuado e o conceito de índice de acoplamento. Este índice é utilizado na sequência de controle ótimo do \mbox{manipulador}.
A hipótese de que o número de juntas ativas seja maior que o número de
passivas  $(n_{a} > n_{p})$  permite o controle ótimo das juntas passivas, uma vez que na etapa de controle destas há mais entradas (torques nos atuadores das juntas ativas), que elementos a controlar (posição das juntas passivas). }
%  Primeiro Parametro: Nacional ou Internacional -- N/I
%  Segundo parametro: Ostensivo, Reservado, Confidencial ou Secreto -- O/R/C/S
\FRDitaOpcoes{N}{O}
% Cria o formulario
\itaFRD

\end{document}
% Fim do Documento. O massacre acabou!!! :-)
